%%%%%%%%%%%%%%%%%%%%%%%%%%%%%%%%%%%%%%%%%%%%%%%%%%%%%%%%%%%%%%%%%%%%%%%%%%%%%%% 
%% SPDX-FileCopyrightText: 2026 CERN (home.cern) 
%%
%% SPDX-License-Identifier: CC-BY-SA-4.0+
%%%%%%%%%%%%%%%%%%%%%%%%%%%%%%%%%%%%%%%%%%%%%%%%%%%%%%%%%%%%%%%%%%%%%%%%%%%%%%% 

% This file is both the standalone tool manual and an appendix fragment
% included by the WRPC user manual.

\newif\ifwrpctoolstandalone
\ifcsname ver@article.cls\endcsname
  \wrpctoolstandalonefalse
\else
  \wrpctoolstandalonetrue
\documentclass[a4paper,12pt]{article}
\usepackage[english]{babel}
\usepackage{fullpage}
\usepackage[T1]{fontenc}
\usepackage{lmodern}
\usepackage{microtype}
\usepackage{xcolor}
\usepackage{hyperref}
\usepackage{listings}
\usepackage{mdwlist}

\lstset{columns=flexible, upquote=true, frame=single,
basicstyle=\footnotesize\ttfamily, backgroundcolor=\color[gray]{0.95}}

\title{WRPC Host Tool}
\author{White Rabbit Team}

\begin{document}
\maketitle
\fi

\section{WRPC host tool}
\label{appendix:wrpc_tool}

The \texttt{wrpc} host utility is a multi-purpose command-line tool for
accessing a running WRPC instance. It provides one command-line interface for
several access methods and operations which were previously provided by
board-specific utilities. The exact command set depends on the options used
when the utility is compiled.

\subsection{Invocation model}

The general form is:

\begin{lstlisting}
wrpc COMMAND [COMMAND-OPTIONS] [BOARD-OPTIONS] [ARGUMENTS]
\end{lstlisting}

The most useful discovery commands are:

\begin{lstlisting}
wrpc help
wrpc help COMMAND
wrpc board
wrpc board BOARD
wrpc --version
\end{lstlisting}

\subsection{Building the tool and WRPC firmware}

The host utility is built from \texttt{tools/wrpc.c} in the \textit{wrpc-sw}
repository. The repository must be configured before building. A typical build
starts as follows:

\begin{lstlisting}
git clone https://gitlab.com/ohwr/project/wrpc-sw.git
cd wrpc-sw
git submodule update --init
make menuconfig
\end{lstlisting}

Use the configuration menu to select the target and features required by the
WRPC design. For a RISC-V WRPC, select the RISC-V architecture and configure
the platform and RAM size to match the FPGA design. The RISC-V cross-compiler
must be installed and available in the path before compiling firmware. For
example:

\begin{lstlisting}
export PATH=/opt/gnu/riscv-11.2-small/bin:$PATH
\end{lstlisting}

There are separate make targets for the embedded firmware and the host
utilities:

\begin{description}
\item[\texttt{make}] Build the configured WRPC firmware and its generated
  images. The main output is normally \texttt{wrc.elf}; depending on the
  selected architecture, binary and memory-initialization outputs such as
  \texttt{wrc.bin}, \texttt{wrc.bram}, \texttt{wrc.mem}, \texttt{wrc.mif}, or
  \texttt{wrc.vhd} may also be generated.
\item[\texttt{make tools}] Build the host-side utilities, including
  \texttt{tools/wrpc}, \texttt{tools/wrpc-load}, and
  \texttt{tools/wrpc-vuart}. This target also builds the host libraries needed
  by the tools.
\item[\texttt{make -C tools wrpc}] Rebuild only the main \texttt{wrpc}
  executable after the software tree has already been configured and its
  generated headers are available.
\end{description}

For a complete firmware and host-tool build, use:

\begin{lstlisting}
make
make tools
\end{lstlisting}

The current configuration can be changed with \texttt{make menuconfig} and
the build can be repeated with \texttt{make}. To remove normal build outputs,
use \texttt{make clean}; \texttt{make distclean} also removes generated
configuration data and is appropriate when starting a configuration from
scratch. Run \texttt{wrpc help} and \texttt{wrpc board} after building to
check the commands and access methods compiled into the resulting executable.

\subsection{Access methods}

Every operation which accesses WRPC registers first selects a board or access
method. If \texttt{-b} is omitted, the default is \texttt{pci}. The available
methods are compile-time dependent; the authoritative list for a particular
binary is printed by \texttt{wrpc board}.

\subsubsection{Generic PCI access}

The \texttt{pci} method maps a PCI resource file. One of \texttt{-f} or
\texttt{-s} is required, and \texttt{-o} selects the WRPC offset within the
resource:

\begin{lstlisting}
wrpc COMMAND -b pci -f /sys/bus/pci/devices/DOMAIN:BUS:DEV.FUNC/resourceN \\
    [-o OFFSET] [COMMAND-OPTIONS]

wrpc COMMAND -b pci -s [DOMAIN:]BUS:DEV[.FUNC][@BAR] [COMMAND-OPTIONS]
\end{lstlisting}

The \texttt{-s} form is converted to the corresponding path under
\texttt{/sys/bus/pci/devices}. Domain, bus, device, function, and BAR are
parsed as hexadecimal values, while the resource-file offset is parsed as a
numeric C-style value.

\subsubsection{SPEC access}

The \texttt{spec} method uses the SPEC PCI resource convention and applies the
WRPC register offset \texttt{0x1000}:

\begin{lstlisting}
wrpc COMMAND -b spec -f RESOURCE-FILE [COMMAND-OPTIONS]
wrpc COMMAND -b spec -s [DOMAIN:]BUS:DEV[.FUNC] [COMMAND-OPTIONS]
\end{lstlisting}

\subsubsection{Host memory access}

The \texttt{host} method maps an AXI-connected WRPC through a host memory
device. The base address is required. The default mapping file is
\texttt{/dev/mem}; \texttt{-f} can select another device such as a UIO
resource:

\begin{lstlisting}
wrpc COMMAND -b host -b BASE [COMMAND-OPTIONS]
wrpc COMMAND -b host -b BASE -f /dev/mem [COMMAND-OPTIONS]
\end{lstlisting}

Here the first \texttt{-b} selects the access method and the second specifies
the mapped base address. The base address is parsed with base detection, so
values such as \texttt{0xA0000000} may be used.

Other access methods, like \texttt{wren-vme}, are available only when their support
was enabled at build time. Use \texttt{wrpc board BOARD} for the options of an
enabled method.

\subsection{Firmware loading and memory operations}

\texttt{load} operates on the WRPC embedded soft-core CPU memory. It does not program
the FPGA configuration bitstream. FPGA programming remains the responsibility
of the board-specific programming tool or boot flow.

\begin{lstlisting}
wrpc load [-c wrpc-v5|wrpc-v4] BOARD-OPTIONS FILENAME
\end{lstlisting}

By default, the command resets the WRPC CPU, writes the image, and releases
the CPU from reset. The supported options are:

\begin{description}
\item[\texttt{-c CORE}] Select the WRPC memory interface, either
  \texttt{wrpc-v5} for the RISC-V design or \texttt{wrpc-v4} for the older
  interface.
\item[\texttt{-v}] Print each word written.
\item[\texttt{-k}] Read back each word and report verification errors.
\item[\texttt{-r}] Do not reset or restart the CPU.
\item[\texttt{-d}] Dump the first \texttt{0x200} bytes instead of loading an
  image.
\item[\texttt{-s}] Save the WRPC memory to the named file. The saved image is
  128 KiB.
\end{description}

Raw binary images are written sequentially from address zero. 32-bit
little-endian RISC-V ELF images are also accepted by the v5 loader; loadable
segments are written to their ELF virtual addresses. The image must be
word-aligned after the loader's four-byte padding rules are applied.

Examples:

\begin{lstlisting}
wrpc load -b pci -s 0000:20:00.0@0 wrc.bin
wrpc load -c wrpc-v5 -b pci -s 0000:20:00.0@0 -k wrc.elf
wrpc load -b pci -s 0000:20:00.0@0 -d
wrpc load -b pci -s 0000:20:00.0@0 -s saved-memory.bin
\end{lstlisting}

For the Kria KR260 reference design, the WRPC is mapped through the host
access method at \texttt{0x80000000}. After copying the rebuilt
\texttt{wrc.elf} to the board, load it without rebuilding the FPGA project:

\begin{lstlisting}
sudo wrpc-sw/tools/wrpc load -b host -b 0x80000000 wrc.elf
\end{lstlisting}

The KR260 reference design uses a 192~kB (196608-byte) WRPC memory. Configure
the WRPC firmware for the RISC-V target, generic 16-bit platform, and this RAM
size before building the firmware.

\subsection{Virtual UART and WRPC shell}

The \texttt{vuart} command accesses the WRPC virtual UART mapped in the core.
With no mode option it provides an interactive terminal. Interactive mode uses
a terminal configuration suitable for the virtual UART and exits when
\texttt{Ctrl-A} is pressed.

\begin{lstlisting}
wrpc vuart BOARD-OPTIONS [-k] [-c COMMAND] [-r] [-t TIMEOUT] [-d]
\end{lstlisting}

\begin{description}
\item[\texttt{-c COMMAND}] Send one WRPC shell command and print its response
  until the \texttt{wrc\#} prompt is detected.
\item[\texttt{-r}] Read-only mode. Print incoming VUART data without reading
  standard input; useful for scripts and logging.
\item[\texttt{-t SECONDS}] Stop terminal or read-only mode after the given
  number of seconds.
\item[\texttt{-k}] Keep the current terminal settings instead of switching
  to raw VUART settings.
\item[\texttt{-d}] Print VUART register activity for debugging.
\end{description}

The \texttt{-c}, \texttt{-r}, and \texttt{-d} modes are mutually exclusive.

Examples:

\begin{lstlisting}
wrpc vuart -b pci -s 0000:20:00.0@0
wrpc vuart -b pci -s 0000:20:00.0@0 -c "ptp status"
wrpc vuart -b pci -s 0000:20:00.0@0 -r -t 10
\end{lstlisting}

On the Kria KR260, connect to the same host mapping used for firmware loading:

\begin{lstlisting}
sudo wrpc-sw/tools/wrpc vuart -b host -b 0x80000000
\end{lstlisting}

This opens the WRPC shell through the virtual UART after the FPGA gateway and
WRPC software have been loaded.

\subsection{Information and diagnostics}

The following commands provide basic WRPC inspection:

\begin{description}
\item[\texttt{info}] Display WRPC hardware and software identification and
  check that the selected mapping is usable.
\item[\texttt{mac}] Read and display the WRPC endpoint MAC address.
\item[\texttt{wdiags}] Read the WR diagnostics snapshot, including servo,
  port, PTP, auxiliary-clock, frame-counter, time, delay, offset, phase, and
  temperature information.
\item[\texttt{aux-logger}] Print the AUX0 diagnostic value in a form suitable
  for logging.
\item[\texttt{rd}] Read one word from the local bus. The command supports an
  address option \texttt{-a ADDRESS}; its default is \texttt{0xff040000}. It
  is primarily useful on access methods which expose a local-bus mapping.
\end{description}

\subsection{SoftPLL tools}

\texttt{spll-recorder} captures real-time SoftPLL traces, such as error
values, DAC drive values, and events. Its documented options are:

\begin{description}
\item[\texttt{-u USAMP}] Set the undersampling factor when supported by the
  selected SoftPLL data source.
\item[\texttt{-d}] Enable debug output.
\item[\texttt{-b}] Produce binary output.
\end{description}

The output can be passed to \texttt{spll-display}, which displays recorded
SoftPLL data. The exact recorder/display options should be checked with
\texttt{wrpc help spll-recorder} and \texttt{wrpc help spll-display}, because
the available data sources depend on the build configuration.

\subsection{Debugging and RPU commands}

\texttt{gdbserver} exposes the WRPC RISC-V CPU through a GDB remote server:

\begin{lstlisting}
wrpc gdbserver BOARD-OPTIONS [-p PORT] [-v] [-t] [-k] [-r]
\end{lstlisting}

The options select the TCP listen port, verbose output, terminal support, and
whether the server keeps accepting connections. The \texttt{-r} option resets
the WRPC CPU before starting the server.

The Zynq UltraScale+ RPU commands provide corresponding operations for the
Cortex-R5 processors:

\begin{description}
\item[\texttt{load-rpu0}, \texttt{load-rpu1}] Load an ARM ELF image into the
  selected RPU and restart it.
\item[\texttt{gdbserver-rpu0}, \texttt{gdbserver-rpu1}] Start a GDB server
  for the selected RPU.
\item[\texttt{checkdbg-rpu0}, \texttt{checkdbg-rpu1}] Inspect and initialize
  the selected RPU debug port.
\item[\texttt{clear-tcm-rpu0}, \texttt{clear-tcm-rpu1}] Clear the selected
  RPU ATCM and BTCM memories.
\item[\texttt{rpu0}, \texttt{rpu1}] Run RPU subcommands. The subcommands
  include \texttt{dfault}, \texttt{tcm}, \texttt{mpu},
  \texttt{rd ADDR [LEN]}, and \texttt{wr ADDR VALUE}.
\item[\texttt{zynqmp-rpu}] Display ZynqMP RPU status.
\end{description}

The RPU commands require a Zynq UltraScale+ access method and are not
available or meaningful on ordinary PCI WRPC deployments.

The exact command set depends on build-time configuration, including the
selected WRPC architecture, target, and optional board support. Therefore the
binary's own help output is the final authority:

\begin{lstlisting}
wrpc help
wrpc help load vuart
wrpc board
\end{lstlisting}

The separate programs \texttt{wrpc-load} and \texttt{wrpc-vuart} are legacy
standalone tools built from their own source files. The \texttt{wrpc}
dispatcher is the generic tool and should be preferred when it is available.

\ifwrpctoolstandalone
\end{document}
\fi